\documentclass[11pt,twoside,openright]{book}

\usepackage[utf8]{inputenc}
\usepackage[T1]{fontenc}
\usepackage{lmodern}
\usepackage{microtype}
\usepackage{enumitem}
\usepackage{hyperref}
\usepackage[
  paperwidth=170mm,
  paperheight=240mm,
  inner=24mm,
  outer=18mm,
  top=22mm,
  bottom=24mm,
  headsep=7mm
]{geometry}

\hypersetup{
  colorlinks=true,
  linkcolor=black,
  citecolor=black,
  urlcolor=blue,
  pdfauthor={Carlo Perassi},
  pdftitle={Book structure proposal}
}

\setlist[itemize]{leftmargin=2em,itemsep=0.25em,topsep=0.4em}
\setlist[enumerate]{leftmargin=2em,itemsep=0.25em,topsep=0.4em}

\title{Book Structure Proposal}
\author{Carlo Perassi}
\date{\today}

\begin{document}

\frontmatter
\maketitle
\tableofcontents

\chapter{Editorial premise}

This file proposes a book-shaped container for the existing draft material. It is
not a final table of contents. Its purpose is to make the physical format,
reader journey, and merge points explicit before the individual essays are
rewritten into a single manuscript.

\section*{Format}

\begin{itemize}
  \item Book class, two-sided layout, chapters opening on right-hand pages.
  \item Trim size: 170mm by 240mm.
  \item Wider inner margin for binding and a compact outer margin for notes.
  \item Plain LaTeX dependencies so individual drafts can be imported gradually.
\end{itemize}

\section*{Working editorial rule}

Each existing document should enter the book only after it has a clear role:
foundation, case study, mathematical appendix, or speculative essay. The book
should avoid becoming a repository dump; the structure should create a sequence
of questions and answers.

\mainmatter

\part{Foundations}

\chapter{The problem space}

Opening chapter for the central question, target reader, and intellectual
stakes. This should define the book's recurring vocabulary and explain why the
chapters belong together.

\section{Candidate functions}

\begin{itemize}
  \item State the core thesis in a form that can survive across technical and
        philosophical chapters.
  \item Introduce the reader to the recurring method: precise models, strategic
        consequences, and cultural interpretation.
  \item Mark what the book intentionally does not try to settle.
\end{itemize}

\chapter{Conceptual primitives}

A short catalogue of indispensable concepts. This can absorb or adapt material
from \texttt{ten\_indispensable.tex}, but should be rewritten as connective
tissue rather than kept as a standalone list.

\part{Mathematical and computational models}

\chapter{Limits, branches, and continuation}

Candidate home for material from \texttt{limit\_complex.tex}. The chapter should
motivate the mathematics before formal development, then explain why the model
matters for the broader argument.

\chapter{Signals and transient spherical models}

Candidate home for \texttt{two\_transient\_spherical\_signal\_models.tex}. This
chapter can work as a technical case study: two models, their assumptions, and
what their comparison teaches.

\chapter{Games native to computation}

Candidate home for \texttt{ai-native-games-v01.tex} and related game material.
The chapter should connect mathematical substrate design to strategic systems,
rather than only presenting a taxonomy.

\part{Culture, institutions, and interpretation}

\chapter[Strategic challenges]{Strategic challenges and civilizational pressure}

Candidate home for \texttt{seven\_strategic\_challenges\_europe.tex}. The
chapter should be framed as an applied analysis that uses the book's primitives,
not as a separate geopolitical essay.

\chapter[Demography and integration]{Demography, integration, and liberal self-preservation}

Candidate home for \texttt{western\_death.tex}. This section likely needs the
strongest editorial framing, because it will set the tone for how controversial
claims are argued, bounded, and evidenced.

\chapter[Religion and explanatory discipline]{Religion, atheism, and explanatory discipline}

Candidate home for \texttt{study\_sheets\_on\_scientific\_atheism.tex}. The
chapter can be organized around what counts as a good explanatory replacement,
not merely around critique.

\part{Synthesis}

\chapter{What world do they want?}

Candidate home for \texttt{what\_world\_do\_they\_want.tex}. This chapter can
serve as the synthetic payoff, asking how explicit models of agency, value, and
constraint change the interpretation of political or technological programs.

\chapter{After the models}

Closing chapter. It should return to the reader's practical question: how to
think better after the abstractions have done their work.

\appendix

\chapter{Appendix: Source draft inventory}

\begin{itemize}
  \item \texttt{ai-native-games-v01.tex}: computational games and strategic
        substrates.
  \item \texttt{limit\_complex.tex}: complex-analysis note on exponent limits.
  \item \texttt{two\_transient\_spherical\_signal\_models.tex}: mathematical
        comparison of transient signal models.
  \item \texttt{seven\_strategic\_challenges\_europe.tex}: strategic pressures
        facing Europe.
  \item \texttt{western\_death.tex}: demography and liberal-democratic
        preservation.
  \item \texttt{study\_sheets\_on\_scientific\_atheism.tex}: structured critique
        and explanatory alternatives.
\end{itemize}

\chapter{Appendix: Merge checklist}

\begin{enumerate}
  \item Decide whether each draft becomes a chapter, section, appendix, or source
        note.
  \item Normalize title capitalization and theorem/list styles across imported
        material.
  \item Replace standalone introductions with chapter-level transitions.
  \item Add a bibliography strategy before importing claims that need citations.
  \item Build the full book after each imported chapter to catch package
        conflicts early.
\end{enumerate}

\backmatter

\chapter{Notes for next revision}

\begin{itemize}
  \item Decide the actual book title and subtitle.
  \item Choose whether the manuscript should be mainly technical, essayistic, or
        hybrid.
  \item Identify which drafts should be excluded from the first integrated
        version.
\end{itemize}

\end{document}
